% GENERATED FILE. Edit canonical lessons, then run npm run generate:books.
% canonical-source-hash: fnv1a64:34afb5795acc2445
% canonical-lessons: MR-C121-sakne, MR-C121-jamne, MR-C121-kausalya, MR-C121-ksamata, MR-C121-adhikar

\chapter{To be able to, To manage to, Skill, Ability, Right}
\label{ch:mr-to-be-able-to-to-manage-to-skill-ability-right}
\hlchaptermodality{121}

\begin{chapteropening}
\textbf{By the end of this chapter:} \emph{I can say to be able to, to manage to, a skill, ability and a right.}

Complete the last lesson of chapter 121: i can say to be able to, to manage to, a skill, ability and a right.
\end{chapteropening}

\section[śakṇe]{\mr{शकणे} (śakṇe) — to be able to, can}

\label{lesson:MR-C121-sakne}



\begin{quote}
Before the new one: say the Marathi for a country, then the Marathi for a story.

Say something you can already do: \textbf{\mr{मी} \mr{मराठी} \mr{बोलतो}}, with no separate word for \emph{am}.
\end{quote}

\subsection*{You'll want to know: \mr{शकणे}}
\textbf{\mr{शकणे}} — \emph{śakṇe} — \textquotedblleft{}to be able to, can\textquotedblright{}.

\begin{grammarlens}[title={What you've built}]
One more word for this chapter.
\end{grammarlens}

\subsection*{Your turn}
\begin{itemize}
\raggedright
  \item \emph{Say it:} \emph{śakṇe}
  \item \emph{Say it:} \emph{śakṇe}, once more
  \item \emph{Say it:} say \emph{śakṇe}
  \item \emph{Recall:} say the Marathi for a country, then the Marathi for a story, then say \emph{śakṇe} again
\end{itemize}

\begin{tcolorbox}[breakable,colback=teal!4,colframe=teal!35!black,title={Before you move on}]
What does \mr{शकणे} mean? (\textquotedblleft{}To be able to, can\textquotedblright{}.) Say it once more.
\end{tcolorbox}
\section[jamṇe]{\mr{जमणे} (jamṇe) — to manage to, to come off}

\label{lesson:MR-C121-jamne}



\begin{quote}
Before the new one: say the Marathi for a story, then the Marathi for to be able to.
\end{quote}

\subsection*{You'll want to know: \mr{जमणे}}
\textbf{\mr{जमणे}} — \emph{jamṇe} — \textquotedblleft{}to manage to, to come off\textquotedblright{}.

\begin{grammarlens}[title={What you've built}]
One more word for this chapter.
\end{grammarlens}

\subsection*{Your turn}
\begin{itemize}
\raggedright
  \item \emph{Say it:} \emph{jamṇe}
  \item \emph{Say it:} \emph{jamṇe}, once more
  \item \emph{Say it:} say \emph{jamṇe}
  \item \emph{Recall:} say the Marathi for a story, then the Marathi for to be able to, then say \emph{jamṇe} again
\end{itemize}

\begin{tcolorbox}[breakable,colback=teal!4,colframe=teal!35!black,title={Before you move on}]
What does \mr{जमणे} mean? (\textquotedblleft{}To manage to, to come off\textquotedblright{}.) Say it once more.
\end{tcolorbox}
\section[kauśalya]{\mr{कौशल्य} (kauśalya) — a skill}

\label{lesson:MR-C121-kausalya}



\begin{quote}
Before the new one: say the Marathi for to be able to, then the Marathi for to manage to.

Many skills are a noun plus \textbf{\mr{करणे}}: say \textbf{\mr{काम} \mr{करणे}}, to work.
\end{quote}

\subsection*{You'll want to know: \mr{कौशल्य}}
\textbf{\mr{कौशल्य}} — \emph{kauśalya} — \textquotedblleft{}a skill\textquotedblright{}.

\begin{grammarlens}[title={What you've built}]
One more word for this chapter.
\end{grammarlens}

\subsection*{Your turn}
\begin{itemize}
\raggedright
  \item \emph{Say it:} \emph{kauśalya}
  \item \emph{Say it:} \emph{kauśalya}, once more
  \item \emph{Say it:} say \emph{kauśalya}
  \item \emph{Recall:} say the Marathi for to be able to, then the Marathi for to manage to, then say \emph{kauśalya} again
\end{itemize}

\begin{tcolorbox}[breakable,colback=teal!4,colframe=teal!35!black,title={Before you move on}]
What does \mr{कौशल्य} mean? (\textquotedblleft{}A skill\textquotedblright{}.) Say it once more.
\end{tcolorbox}
\section[kṣamatā]{\mr{क्षमता} (kṣamatā) — ability, capacity}

\label{lesson:MR-C121-ksamata}



\begin{quote}
Before the new one: say the Marathi for to manage to, then the Marathi for a skill.
\end{quote}

\subsection*{You'll want to know: \mr{क्षमता}}
\textbf{\mr{क्षमता}} — \emph{kṣamatā} — \textquotedblleft{}ability, capacity\textquotedblright{}.

\begin{grammarlens}[title={What you've built}]
One more word for this chapter.
\end{grammarlens}

\subsection*{Your turn}
\begin{itemize}
\raggedright
  \item \emph{Say it:} \emph{kṣamatā}
  \item \emph{Say it:} \emph{kṣamatā}, once more
  \item \emph{Say it:} say \emph{kṣamatā}
  \item \emph{Recall:} say the Marathi for to manage to, then the Marathi for a skill, then say \emph{kṣamatā} again
\end{itemize}

\begin{tcolorbox}[breakable,colback=teal!4,colframe=teal!35!black,title={Before you move on}]
What does \mr{क्षमता} mean? (\textquotedblleft{}Ability, capacity\textquotedblright{}.) Say it once more.
\end{tcolorbox}
\section[adhikār]{\mr{अधिकार} (adhikār) — a right, authority}

\label{lesson:MR-C121-adhikar}



\begin{quote}
Before the new one: say the Marathi for a skill, then the Marathi for ability.
\end{quote}

\subsection*{You'll want to know: \mr{अधिकार}}
\textbf{\mr{अधिकार}} — \emph{adhikār} — \textquotedblleft{}a right, authority\textquotedblright{}.

\begin{grammarlens}[title={What you've built}]
One more word for this chapter.
\end{grammarlens}

\subsection*{Your turn}
\begin{itemize}
\raggedright
  \item \emph{Say it:} \emph{adhikār}
  \item \emph{Say it:} \emph{adhikār}, once more
  \item \emph{Say it:} say \emph{adhikār}
  \item \emph{Recall:} say the Marathi for a skill, then the Marathi for ability, then say \emph{adhikār} again
\end{itemize}

\begin{tcolorbox}[breakable,colback=teal!4,colframe=teal!35!black,title={Before you move on}]
What does \mr{अधिकार} mean? (\textquotedblleft{}A right, authority\textquotedblright{}.) Say it once more.
\end{tcolorbox}
